Software is invisible. There is no smoke when a program runs, no exhaust pipe, no physical waste left behind. Yet as this experiment demonstrates, the choice of how to write software, specifically, which programming language to use, produces measurable and statistically significant differences in the carbon emissions. In a world where the ICT sector already accounts for an estimated 2.1--3.9\% of global greenhouse gas emissions \cite{ict_stat}, and where that share is growing with the expansion of cloud computing and artificial intelligence, the environmental cost of software design decisions is no longer a niche concern. It is a mainstream one. 

This paper presented a controlled empirical study measuring the carbon footprint of six programming languages, across six algorithms and three input sizes on a single hardware platform. Energy was estimated using the Intel RAPL accesses through LibreHardwareMonitor, converted to gCO$_2$e using real-time grid carbon intensity from the Electricity Maps API, and analyzed using Welch's $t$-tests with Bonferroni correction across all 15 pairwise language comparisons. The dataset comprised 12,420 measurements across 108 experimental cells, with 115 replications per cell.

The results are statistically robust and practically significant. \textbf{14 of 15 language pairs are distinguishable in carbon emissions} at the Bonferroni-corrected threshold of $\alpha' = 0.0033$. The sole exception is C versus Rust ($t = -0.313$, $p = 0.754$, $d = 0.017$) --- a decisive non-result confirming that Rust achieves C-equivalent carbon efficiency despite its compile-time memory safety guarantees. The language landscape divides into four tiers: C and Rust at the efficient end ($1.0\times$), Java at a modest overhead ($1.8\times$), Go and JavaScript at moderate overhead ($3.0\times$--$4.3\times$), and Python at $213\times$ the carbon footprint of C. Python's overhead is not uniform --- it ranges from $1.5\times$ on binary search to $255.6\times$ on matrix multiplication, growing proportionally with algorithmic complexity and confirming that interpreter dispatch overhead compounds with every loop iteration.

The C--Rust equivalence result carries particular significance for the software industry. Rust is increasingly adopted as a safer replacement for C in operating systems, embedded software, and infrastructure components~\cite{google2025rustAndroid}~\cite{linux2026rust}. Our data confirms that this transition carries no energy or carbon cost, a finding that should encourage, not discourage, the adoption of memory-safe systems programming languages.

Beyond the technical findings, this study is a reminder that software engineering decisions exist within a broader environmental context. A developer choosing a programming language is not only making a decision about performance, productivity, and ecosystem, they are, knowingly or not, making a decision about energy and carbon. When that decision is made at scale, across millions of daily executions, across entire cloud deployments, across the global software industry, the aggregate environmental consequence becomes material. The 1,170~kg CO$_2$e per year estimated for a high-throughput Python service versus its C equivalent is not an abstraction. It is the same order of magnitude as driving a car for several thousand kilometers.
 
Making software carbon-visible is one step toward making it carbon-conscious. This study provides a replicable methodology, a concrete empirical baseline, and a four-tier conceptual framework for thinking about the energy cost of language choice. We hope it contributes to a broader culture of sustainability awareness in software development. One where environmental responsibility is considered alongside performance, correctness, and cost as a first-class engineering concern.
